\section{The transformation laws and the actual complex uniformization}
\label{sec:normalized-frame-laws-v051}

\status{NFL1--NFL22 prove the remaining three fields of
\texttt{FramesOK}, the original \texttt{CoverTri} and \texttt{FibTri},
and the full \texttt{IsUnf} and principal/cusp return on the same family
constructed in NFR13.  The supplied \texttt{Moduli} record and the exact
\texttt{SecBaseChangeSt} and \texttt{SecTwistSt} laws remain premises.
No whole selected RH/GRH obligation is declared closed.  All source
identifications refer to commit
\texttt{7cbf7a0c32fca1badf05931879a37c3f8a2b03b4}.}

We retain all the original data and conventions of NFR1--NFR13:
$E,H,D_w,l,X_o,c$, $\sigma,g,e,F,r,M$, with \texttt{GoodEmb},
\texttt{ModelBase} and \texttt{NegXi}.  Thus $F$ is the actual
\texttt{FamilyUnif} of the supplied moduli problem, and $M$ its model
on the prescribed Hermitian space.  Fix a component $k$ temporarily.
Write $\Lambda=F.D_i(k).\mathrm{La}$, $G=M.X_k.G_m$,
$\xi_0=M.X_k.\xi_0$, $\phi=M.\phi_k$, and

\[
 E_\Lambda(v,w)=\operatorname{Tr}_{E/\mathbb Q}
       \bigl(-\xi_0\operatorname{herm}_G(\phi v,\phi w)\bigr).
 \tag{NFL.1}
\]

Hermitian forms are linear in the first argument.  The first $q$
columns of a frame have positive sign, and the remaining $r-q$
columns negative sign.  Let $\mathcal D$ be the actual open flag
domain, $T$ the \texttt{ThetaData} of NFR4, and
$b=T.\mathrm{co}^{-1}x$ its linear coordinates.  Its period map,
normalizer, Siegel matrix and rigidified local frames are denoted
$P_b,h_b,\tau_b,fr_j(b,z)$.  Its two fixed half-characteristics are
$\alpha,\beta\in\{0,1/2\}^g$.  The actual normalized point
$u_k(x)$ has section

\[
 s_b(z)=\kappa_b\theta[\alpha,\beta](\tau_b,h_bz)fr_j(b,z),
 \qquad \kappa_b^2=C\det h_b,
 \quad C=\operatorname{normConst}(\mathrm{ty},\xi_0)\ne0.
 \tag{NFL.2}
\]

The square root is chosen separately at each point.  No continuity
of $b\mapsto\kappa_b$ is assumed.  NFR12 already proves every field
of \texttt{FrameOK} for these exact points, including its algebraic
section and mean-one condition.

Here are the retained laws with their actual scope.  The
\texttt{SecBaseChangeSt} law starts with a locally Noetherian scheme
$S$, a polarized abelian variety $P/S$, and $n\ge1$.  Near each
point it supplies $n^{\dim P}$ sections whose pullbacks form a basis
after \emph{every} affine base change $\operatorname{Spec}R\to S$;
$R$ need not be reduced, Noetherian, or flat over the original base.
We use the $n=1$ instance of the actual principally polarized object
attached to $X_o$.  The \texttt{SecTwistSt} law, for every $S,P,q\ge1$
and every local ring $R$ with $t:\operatorname{Spec}R\to S$ and
$c\in P(R)$, supplies a basis of $q^{\dim P}$ sections of
$L^{q-1}\otimes t_c^*L$ and lifts every residue-field section.
The arguments below use its original complex-fiber instances.
Neither law is replaced by an assertion only for flat extensions.

We use HUF21 for the actual holomorphic family, its exponential,
lattice fibers and theta frames; PUN11 for its models; and TDR1--TDR13
for literal dual-number eigenframes, section generators, universal
frame rings and descent.  Their already stated classical foundations
are retained.  The additional analytic ingredients below are ordinary
Gaussian integration, Poisson summation for Schwartz functions,
Cauchy's theorem, Liouville's theorem and local analytic coordinates
on smooth complex algebraic varieties.  The needed theta multiplier
is proved, rather than made an extra input.

\subsection{The square of the theta multiplier}

% BEGIN NFL1
\paragraph{NFL1: the finite quadratic sum with its exact sign.}
\label{claim:NFL1}
Let $\gamma=\left(\begin{smallmatrix}A&B\\ C&D\end{smallmatrix}\right)
\in\operatorname{Sp}_{2g}(\mathbb Z)$, with $B,C$ even and $\det D$
odd.  On $Q=\mathbb Z^g/D\mathbb Z^g$ put
\[
 S(B,D)=\sum_{v\in Q}\exp\bigl(\pi i v^tBD^{-1}v\bigr).
\]
Then the summand is well defined and
\[
 S(B,D)^2=\chi_4(\det D)\det D,
 \qquad
 \chi_4(d)=\begin{cases}1&d\equiv1\pmod4,\\-1&d\equiv3\pmod4.
 \end{cases}
 \tag{NFL.3}
\]
The assertion includes $g=0$, with all determinants and sums equal
to one.

\begin{proof}
Symplecticity gives $BD^{-1}=(BD^{-1})^t$ and $D^tB$ symmetric.
Replacing $v$ by $v+Dn$ changes $v^tBD^{-1}v$ by
$2n^tB^tv+n^tD^tBn$, an even integer.  The resulting quadratic
function $q(v)=\tfrac12v^tBD^{-1}v\pmod{\mathbb Z}$ has associated
pairing $b(v,w)=v^tBD^{-1}w\pmod{\mathbb Z}$.
If this pairing kills $w$, then $BD^{-1}w$ is integral.  Both
$B(D^{-1}w)$ and $D(D^{-1}w)$ are integral, so
$(A^tD-C^tB)D^{-1}w=D^{-1}w$ is integral.  Hence $w=0$ in $Q$.
Thus $q$ is nondegenerate and $|Q|=|\det D|$ is odd.

For completeness, an odd finite abelian group with a nondegenerate
symmetric pairing into $\mathbb Q/\mathbb Z$ splits orthogonally
into cyclic prime-power groups with nondegenerate self-pairing.
First split its primary parts.  On a $p$-part choose an element of
largest order $p^e$.  Nondegeneracy pairs some two elements with
denominator $p^e$.  If neither has self-pairing with that denominator,
their sum does: its self-pairing contains twice their pairing, and
$p$ is odd.  The cyclic subgroup it generates is nondegenerate.
Subtracting the appropriate multiple of that generator projects
every element onto its orthogonal complement.  Induction proves
the assertion.  Since $2$ is invertible on an odd group, $q$ is
determined by $b(v,v)/2$; thus the Gauss sum factors across this
orthogonal decomposition.

On a cyclic summand it is
$G(a,p^e)=\sum_{x\bmod p^e}\exp(2\pi i ax^2/p^e)$ with $p\nmid a$.
For $e\ge2$, write $x=y+tp^{e-1}$ and sum over $t\bmod p$.
The inner sum vanishes unless $p\mid y$, and gives
$G(a,p^e)=pG(a,p^{e-2})$ otherwise.  For $e=1$, if $\chi$ is the
quadratic character of $\mathbb F_p$, counting square roots gives
$G(a,p)=\chi(a)\tau$, where
$\tau=\sum_x\chi(x)e^{2\pi ix/p}$.
Character orthogonality gives $|\tau|^2=p$: expand the square,
put $x=ty$ for $y\ne0$, and use
$\sum_{y\ne0}e^{2\pi i(t-1)y/p}=p-1$ for $t=1$ and $-1$ otherwise.
Also $\overline\tau=\chi(-1)\tau$.  Therefore
$G(a,p)^2=\chi(-1)p$, and the recurrence yields
$G(a,p^e)^2=\chi_4(p)^e p^e$.
Multiplication over all summands gives
$S^2=\chi_4(|\det D|)|\det D|$.
This equals the right side of (NFL.3) also when $\det D<0$,
because $\chi_4(-1)=-1$.
\end{proof}
% END NFL1

% BEGIN NFL2
\paragraph{NFL2: the exact theta-square transformation.}
\label{claim:NFL2}
For a Siegel matrix $\tau$, put
\[
 J=C\tau+D,\qquad \tau'=(A\tau+B)J^{-1},
 \qquad Q_\gamma(z)=z^tJ^{-1}Cz.
\]
If $B,C$ are divisible by four, then, for integer vectors $k,k',m,m'$
with $m\equiv A^tk\pmod2$ and $m'\equiv D^tk'\pmod2$,
\[
 \begin{split}
 \theta[k/2,k'/2](\tau',J^{-t}z)^2
   ={}&\chi_4(\det D)\det J\,e^{2\pi iQ_\gamma(z)}\\
     &{}\cdot\theta[m/2,m'/2](\tau,z)^2.
 \end{split}
 \tag{NFL.4}
\]
In particular this proves the original zero-characteristic
\texttt{theta\_sq\_transform\_level\_eight} law.

\begin{proof}
We give the multiplier computation.  For $\operatorname{Im}T>0$ let
$\mathcal G(T)=\int_{\mathbb R^g}e^{\pi it^tTt}\,dt$.
Real Gaussian integration and holomorphic continuation on the
connected Siegel space give
$\mathcal G(T)^2\det T=i^g$.
The Fourier transform of $e^{\pi it^tTt+2\pi it^tz}$ is
$\mathcal G(T)e^{-\pi i(z-m)^tT^{-1}(z-m)}$ at $m\in\mathbb Z^g$.
Poisson summation therefore gives
\[
 \theta[0,0](\tau,z)=\mathcal G(\tau)
 e^{-\pi iz^t\tau^{-1}z}
 \theta[0,0](-\tau^{-1},\tau^{-1}z).
 \tag{NFL.5}
\]
The Gaussian decay permits the Fourier transform, summation and
holomorphic continuation just used; derivatives on compact subsets
are bounded by a polynomial times a decaying real Gaussian.

Since $AD^t-BC^t=I$ and $B,C$ are even, $D$ is invertible modulo
two, so $\det D$ is odd and nonzero.  Set
$N=\tau^{-1}+D^{-1}C$.  The real symmetric matrix $D^{-1}C$
does not change the negative definite imaginary part of $\tau^{-1}$;
thus $N^{-1}$ is Siegel.  Block multiplication gives
\[
 J=DN\tau,\quad
 \tau'=D^{-t}N^{-1}D^{-1}+BD^{-1},\quad
 J^{-t}z=D^{-t}N^{-1}\tau^{-1}z.
 \tag{NFL.6}
\]
Decompose the series defining the theta function at $\tau'$ into
the cosets $v+D\mathbb Z^g$.  Complete the square in each coset
and use (NFL.5).  The finite Fourier sum that occurs is
\[
 \sum_{v\bmod D\mathbb Z^g}
 e^{\pi iv^tBD^{-1}v+2\pi im^tD^{-1}v}
   =S(B,D)e^{\pi im^tD^{-1}Cm}.
 \tag{NFL.7}
\]
Here replacing $v$ by $v-Cm$ proves the identity: use
$BC^t=AD^t-I$, the symmetries of $BD^{-1}$ and $D^{-1}C$,
and that $m^tA^tCm$ is even.  All remaining terms are integral
multiples of $2\pi i$.  Substitution into the coset sum yields
\[
 \theta[0,0](\tau',J^{-t}z)e^{-\pi iQ_\gamma(z)}
   =\lambda\theta[0,0](\tau,z),\qquad
 \lambda=S(B,D)\frac{\mathcal G(N^{-1})}{\mathcal G(\tau)}.
 \tag{NFL.8}
\]
Equations (NFL.3), (NFL.6) and the squared Gaussian formula give
$\lambda^2=\chi_4(\det D)\det J$.  This computes the sign without
choosing a square root of $\det J$.

To include the characteristics, use the series identity
\[
 \theta[a,b](\tau,z)=
 e^{\pi ia^t\tau a+2\pi ia^t(z+b)}
 \theta[0,0](\tau,z+\tau a+b).
 \tag{NFL.9}
\]
Put $n=A^tk+C^tk'$ and $n'=B^tk+D^tk'$.
Insert (NFL.9) into (NFL.8) with the translated argument and
expand its quadratic exponent.  The terms depending on $z$ give
$e^{\pi iQ_\gamma(z)}$; the remaining characteristic factor before
squaring is $e^{\pi i(k\cdot k'-n\cdot n')/4}$.
Indeed the required shift equality follows by multiplying by
$J^t$ and using $A^tD-C^tB=I$; expansion of the constant
quadratic terms then leaves $(k\cdot k'-n\cdot n')/4$ in the
exponent $\pi i$.  The symplectic identity also gives
\[
 k\cdot k'=(A^tk)\cdot(D^tk')-(B^tk)\cdot(C^tk').
\]
Consequently $k\cdot k'-n\cdot n'$ is divisible by four when
$B,C$ are.  Its squared exponential is one.  Changing an integer
representative of a half-characteristic multiplies theta by a
sign, so leaves its square unchanged.  Since
$n\equiv m$ and $n'\equiv m'$ modulo two, (NFL.4) follows.
\end{proof}
% END NFL2

% BEGIN NFL3
\paragraph{NFL3: a rigidified unit fixes the characteristic and multiplier square.}
\label{claim:NFL3}
Suppose $B,C$ are divisible by eight, $\chi_4(\det D)=1$, and
$\rho:\mathbb C^g\to\mathbb C^*$ is entire with $\rho(0)=1$.
If $\lambda\ne0$ and
\[
 \theta[\alpha,\beta](\tau',J^{-t}z)\rho(z)
      =\lambda\theta[\alpha,\beta](\tau,z),
 \tag{NFL.10}
\]
then $A^t(2\alpha)\equiv2\alpha\pmod2$,
$D^t(2\beta)\equiv2\beta\pmod2$, and $\lambda^2=\det J$.
This holds even when the theta function vanishes at zero.

\begin{proof}
Write $n=2\alpha,n'=2\beta$ and
$r=\lambda/\rho$.  This is an entire nowhere-zero function and
$\theta(\tau',J^{-t}z)=r(z)\theta(\tau,z)$.
For $\ell=\tau m+m'$ set
$p=Dm-Cm'$ and $p'=Am'-Bm$; then
$J^{-t}\ell=\tau'p+p'$.
The translation formula proved in NFR5, applied on the nonzero
set of theta and then extended by continuity, shows that
$s(z)=r(z)e^{-\pi iQ_\gamma(z)}$ satisfies
\[
 s(z+\ell)=(-1)^{L(m,m')}s(z),\quad
 L=n\cdot p'-n'\cdot p-n\cdot m'+n'\cdot m
          -m\cdot m'+p\cdot p'.
 \tag{NFL.11}
\]
One can check the quadratic part directly from
\[
 \begin{split}
 Q_\gamma(z+\ell)-Q_\gamma(z)
   ={}&2m\cdot z+m\cdot\tau m
       -2p\cdot J^{-t}z-p\cdot\tau'p\\
     &{}+m\cdot m'-p\cdot p'.
 \end{split}
\]
It follows that $s^2$ is periodic for the full lattice
$\mathbb Z^g+\tau\mathbb Z^g$.  It is bounded on a compact
fundamental parallelepiped, hence everywhere.  Liouville's theorem
on complex lines makes it constant.  Since $s$ never vanishes,
it too is constant (its image lies in the two square roots of a
nonzero constant).  Thus every $L(m,m')$ is even.
Taking $m=0,m'=e_i$, and then $m=e_i,m'=0$, and using that
$B,C$ are even, gives the two asserted parity conditions.

NFL2 now applies with these unchanged half-characteristics.
Squaring (NFL.10) gives
\[
 [\det J\,e^{2\pi iQ_\gamma(z)}\rho(z)^2-\lambda^2]
       \theta[\alpha,\beta](\tau,z)^2=0.
\]
The theta function is not identically zero by NFR5.  The bracket
therefore vanishes on a nonempty open set and hence everywhere
by the identity theorem.  Evaluating the bracket, not a quotient
of theta functions, at zero gives $\lambda^2=\det J$.
\end{proof}
% END NFL3

\subsection{Level structure and comparison of the actual sections}

% BEGIN NFL4
\paragraph{NFL4: the adapted lattice forces the required level-eight sign.}
\label{claim:NFL4}
Let $a:\Lambda\simeq_{\mathcal O_E}\Lambda$ preserve $E_\Lambda$
and satisfy the original two level conditions: it is the identity
modulo $l.M$ on $\overline{l.M_b}\Lambda$, and the identity modulo
$l.A$ on all of $\Lambda$.  In the adapted symplectic basis of
NFR3, write its integral matrix as
$N=(N_{ij})_{1\le i,j\le2}$.  Then
\[
 N_{12}\equiv N_{21}\equiv0\pmod8,\qquad
 \det N_{11}\equiv1\pmod8.
 \tag{NFL.12}
\]
These conclusions also hold for $a^{-1}$.

\begin{proof}
PUN6 upgrades the original half-level congruence to
$(a-I)\Lambda\subset l.M\Lambda$; it uses the complementary
host ideals and the perfect trace pairing, so does not assume
that the lattice is free over $\mathcal O_E$.
Extend $a$ to $V=\Lambda\otimes_{\mathcal O_E}E$.
Its $\mathcal O_E$-linearity and preservation of (NFL.1) imply
preservation of $\operatorname{herm}_G$: replace one variable
by all $s\in E$ and use nondegeneracy of the trace pairing.
Thus $\delta=\det_Ea$ has absolute value one at every complex
embedding.  The endomorphism $(a-I)/l.M$ stabilizes the finite
projective lattice, so its characteristic coefficients are
integral.  Expanding $\det(I+l.M\,T)$ shows
$(\delta-1)/l.M\in\mathcal O_E$.  Since $l.M\ge3$, every
conjugate of this algebraic integer has absolute value at most
$2/l.M<1$.  If it were nonzero, its nonzero integer norm would
have absolute value less than one.  Hence $\delta=1$.

Let $e_0=w_0^4\in\mathcal O_E$ be the host element used in NFR3.
It satisfies $e_0^2\equiv e_0\pmod8$ and acts in that basis by
$\operatorname{diag}(I_g,0)$ modulo eight.
The map $a$ commutes with $e_0$, so its off-diagonal blocks are
zero modulo eight.  Consider $T=e_0a+(1-e_0)I$ on $V$.
Modulo $8\mathcal O_E$, the idempotent $e_0$ splits the base ring
into its two factors.  On those factors $T$ is respectively
$a$ and $I$, whence $\det_ET\equiv1\pmod8$.
This determinant calculation is valid on the finite projective
lattice: localize where it is free, compute the determinant
polynomial there, and glue its integral coefficients.
Taking $E/\mathbb Q$ norms gives
$\det_{\mathbb Q}T\equiv1\pmod8$.
Its integral matrix on $\Lambda/8\Lambda$ is
$\operatorname{diag}(N_{11},I_g)$, so its determinant is
$\det N_{11}$ modulo eight.  This proves (NFL.12).
The inverse preserves the same levels and form, and the same
argument applies to it.
\end{proof}
% END NFL4

% BEGIN NFL5
\paragraph{NFL5: the precise change of period normalization.}
\label{claim:NFL5}
Suppose $P_{b'}(a v)=H_0P_b(v)$ for the lattice isometry in NFL4,
where $H_0\in\operatorname{GL}_g(\mathbb C)$.  Let $N$ be the
matrix of $a^{-1}$ in the adapted basis, and set
\[
 \gamma=\begin{pmatrix}N_{22}^t&N_{12}^t\\N_{21}^t&N_{11}^t\end{pmatrix},
 \quad J=N_{21}^t\tau_b+N_{11}^t.
\]
Then $\gamma$ is integral symplectic, its off-diagonal blocks are
divisible by eight, and $\chi_4(\det\gamma_{22})=1$.  Moreover
\[
 \tau_{b'}=\gamma\tau_b,\qquad
 h_{b'}H_0=J^{-t}h_b,\qquad
 \det h_b=\det J\det h_{b'}\det H_0.
 \tag{NFL.13}
\]

\begin{proof}
Invertibility of $H_0$ follows already from the period relation
if it is initially specified only as a square matrix: $a$ is
surjective and $P_{b'}(\Lambda)$ spans the underlying real
vector space, so $H_0$ is surjective and hence invertible.
Preservation of the integral alternating form gives
$N^t\left(\begin{smallmatrix}0&I\\-I&0\end{smallmatrix}\right)N
=\left(\begin{smallmatrix}0&I\\-I&0\end{smallmatrix}\right)$.
Transposing its four block equations proves symplecticity of
the displayed $\gamma$.  NFL4 gives its congruences and sign.
Normalize $P_b$ to the matrix $[I\ \tau_b]$ in the fixed basis.
The periods of $P_ba^{-1}$ are then
$[A_0\ B_0]$, where
$A_0=N_{11}+\tau_bN_{21}$ and
$B_0=N_{12}+\tau_bN_{22}$.
Its normalizer is $A_0^{-1}h_b$.  Since
$P_ba^{-1}=H_0^{-1}P_{b'}$, this is $h_{b'}H_0$.
Now $A_0^t=J$ and the normalized second block is
$A_0^{-1}B_0=B_0^t(A_0^t)^{-1}$, by symmetry of the normalized
Siegel matrix.  These are precisely (NFL.13).
\end{proof}
% END NFL5

% BEGIN NFL6
\paragraph{NFL6: the determinant identity for rigidified theta sections.}
\label{claim:NFL6}
In NFL5 the two actual base points are equal, and their
exponentials satisfy $\exp_{b'}(H_0z)=\exp_b(z)$.
On every chart of the actual universal line, put
$\Theta_b(z)=\theta[\alpha,\beta](\tau_b,h_bz)fr_j(b,z)$.
Then
\[
 \det h_b\,\Theta_b(z)^2
   =\det h_{b'}\det H_0\,\Theta_{b'}(H_0z)^2.
 \tag{NFL.14}
\]
The assertion compares values in the same rigidified line, not
sections of two merely isomorphic abstract theta bundles.

\begin{proof}
The lattice isometry, its two level conditions, and the equality
of periods meet exactly the \texttt{FamilyUnif} identification
clause reconstructed in HUF21.  It gives equality of the
moduli points and the displayed exponential identity.
On a chart containing $\exp_bz$ form
$\rho_0(z)=fr_j(b',H_0z)/fr_j(b,z)$.
The actual cocycle equations for the two frames make these
quotients agree on overlaps.  Hence they glue to an entire
nowhere-zero function on $\mathbb C^g$.  The rigidification at
the zero section gives $\rho_0(0)=1$.
The two actual algebraic theta sections are nonzero in the same
one-dimensional section space.  This follows from the retained
complex-fiber section law, or directly from NFR5's Fourier
classification for this fiber.  Consequently for one constant
$\lambda\ne0$,
\[
 \theta[\alpha,\beta](\tau_{b'},h_{b'}H_0z)\rho_0(z)
       =\lambda\theta[\alpha,\beta](\tau_b,h_bz).
\]
Replace $z$ by $h_b^{-1}w$.  The resulting unit $\rho(w)$ has
value one at zero, and NFL5 puts this identity in the form of
NFL3.  Thus $\lambda^2=\det J$.
Since $\Theta_{b'}(H_0z)=\lambda\Theta_b(z)$, the determinant
relation in NFL5 proves (NFL.14).  No division by a theta value
at zero, or anywhere on its divisor, has been used.
\end{proof}
% END NFL6

% BEGIN NFL7
\paragraph{NFL7: change of an actual theta point, with its literal section map.}
\label{claim:NFL7}
Let $y,y'$ be the original \texttt{IsThetaPt} at $b,b'$ with
nonzero scalars $\kappa,\kappa'$, and suppose the periods are
related as in NFL5.  Then their base points agree, their eigenframes
satisfy
\[
 u'_j=d\exp_b(H_0^{-1}e_j),
\]
and there is a unique $t_0\in\mathbb C^*$ whose actual section
base-change relation is $s'=t_0s$.  It obeys
\[
 t_0^2\kappa^2\det h_{b'}\det H_0
       =\kappa'^2\det h_b.
 \tag{NFL.15}
\]
Here $s'=t_0s$ means the literal \texttt{SecBC} over the identity:
the canonical map of the two fiber products respects both
projections and pulls the scaled section to the indicated section
on every original line chart.

\begin{proof}
Differentiate $\exp_{b'}H_0=\exp_b$ using the original
dual-number interpretation from HUF1 and TDR1.  This gives the
eigenframe formula, including its prescribed CM type.
Both sections generate the same rank-one section module, so
$s'=t_0s$ for a unique unit.  The canonical fiber-product map
over the equal base point supplies the two projection equations
and chart pullback equations in \texttt{SecBC}; uniqueness of
the generator makes the scalar independent of the chart.
Use $s=\kappa\Theta_b$ and $s'=\kappa'\Theta_{b'}$, square their
relation and substitute (NFL.14).  Choose any $z$ where theta
is nonzero and a chart containing its exponential image.
The frame value is nonzero there, so cancellation gives
(NFL.15).  Such a point exists by NFR5.
\end{proof}
% END NFL7

\subsection{The two exact transformation laws, for every allowed weight}

% BEGIN NFL8
\paragraph{NFL8: the period matrix of an upper triangular frame move.}
\label{claim:NFL8}
Let $x\in\mathcal D$, let $b\in\operatorname{Bor}(\Psi,r)$, and
put $x'=xb$.  There is an invertible matrix $H_0$ with
$P_{x'}=H_0P_x$.  In the original coordinates indexed by
$j\leftrightarrow(\rho,i)$, both $H_0$ and $H_0^{-1}$ have support
only when the two $\rho$'s agree and the row index $i$ is at most
the column index $i'$.  Write $d_j=(H_0^{-1})_{jj}$.
If $t_0^2\det H_0=1$, then, for every original $\operatorname{Dom}(a,n)$,
\[
 t_0^n\prod_jd_j^{a((ej)_1,(ej)_2)}
       =\prod_{\tau\in\Psi}\prod_{s=0}^{r-1}
                       b_{ss}(\tau)^{\lambda_{\tau,s}},
 \tag{NFL.16}
\]
where $n=2m$ and, using the NFR1 relabelling,
\[
 \lambda_{\tau,s}=
 \begin{cases}
 -a(\operatorname{relab}^{-1}\bar\tau,q-1-s)-m,&s<q,\\
 a(\operatorname{relab}^{-1}\tau,s-q)+m,&s\ge q.
 \end{cases}
 \tag{NFL.17}
\]
These are exactly the source's $\operatorname{laU}$ and $bchar$.

\begin{proof}
Use the literal period formulas of NFR1.  For $s<q$, a period
coordinate is $\operatorname{herm}_{G_\tau}(\operatorname{col}_sX_\tau,
\tau(\phi v))$.  Replacing $X$ by $Xb$ multiplies this early
block by $b_{++}^t$.  For $s\ge q$ it is
$(X_\tau^{-1}\tau(\phi v))_s$, so the late block is multiplied
by $(b_{--})^{-1}$.  The other possible block vanishes because
$b$ is upper triangular.  Thus, in unreversed column order,
\[
 H_0=\bigoplus_{\tau\in\Psi}
            \operatorname{diag}(b_{++}^t,b_{--}^{-1}).
 \tag{NFL.18}
\]
The early columns are reversed by the NFR1 indexing, converting
the first transpose into the indicated upper triangular support.
Inverting preserves that support.  Their diagonals are
$b_{ss}^{-1}$ for $s<q$ and $b_{ss}$ for $s\ge q$.
Also $\det H_0=\prod_{\tau,s<q}b_{ss}
\prod_{\tau,s\ge q}b_{ss}^{-1}$.
Since $n=2m$, $t_0^n=(\det H_0)^{-m}$.  Inserting these three
identities gives (NFL.16) one diagonal entry at a time.
The antitone condition in \texttt{Dom} is retained, although this
particular identity only uses that $n$ is even.  No bound on
$a,n$, and no one-dimensional restriction on the eigenblocks,
has been made.
\end{proof}
% END NFL8

% BEGIN NFL9
\paragraph{NFL9: the complete normalized Borel law.}
\label{claim:NFL9}
For every original $k,x,b$ with $x\in\mathcal D$,
$b\in\operatorname{Bor}(\Psi,r)$ and $xb\in\mathcal D$, there is
\emph{one} $t:\operatorname{Ix}(D_w)\to\mathbb C^*$ such that
\[
 \operatorname{FrBor}(t,u_k(x),u_k(xb))
\]
and, simultaneously for every $a,n$ with $\operatorname{Dom}(a,n)$,
\[
 \operatorname{wchar}(\operatorname{ex}(D_w,a,n))(t)
      =\operatorname{bchar}(\operatorname{laU}(a,n,k))(b).
 \tag{NFL.19}
\]
This is the full original \texttt{normUn\_bor} statement.

\begin{proof}
Apply NFL7 to the identity lattice isometry and the matrix $H_0$
of NFL8.  The two points have equal base point, and the new
eigenvector $j$ is the derivative of $H_0^{-1}e_j$.
Its coefficient at old coordinate $j'$ is
$(H_0^{-1})_{j'j}$, zero unless its original embedding is the
same and its original index is at most that of $j$.
Set $t(\operatorname{some}(\rho,i))=(H_0^{-1})_{jj}$ and
$t(\operatorname{none})=t_0$, the actual section scalar from NFL7.
The diagonal entries are units.  Thus the derivative coefficients,
the equal base point and the literal section pullback of NFL7
give every clause of \texttt{FrBor}.
Substitute the two identities $\kappa_x^2=C\det h_x$ and
$\kappa_{xb}^2=C\det h_{xb}$ into (NFL.15).
All three determinants and $C$ are nonzero, so it reduces to
$t_0^2\det H_0=1$.  NFL8 now gives (NFL.19) for every allowed
weight using the same $t$.  In particular its choice is made
before $a,n$ and before any section to be evaluated.
\end{proof}
% END NFL9

% BEGIN NFL10
\paragraph{NFL10: the complete arithmetic invariance law.}
\label{claim:NFL10}
For every $\gamma\in(M.X_k).\Gamma$ and every $x\in\mathcal D$,
with the original domain guard on $\gamma x$, there is one $t$
such that
\[
 \operatorname{FrBor}(t,u_k(x),u_k(\gamma x)),\qquad
 \operatorname{wchar}(\operatorname{ex}(D_w,a,n))(t)=1
 \tag{NFL.20}
\]
for all $\operatorname{Dom}(a,n)$.  This is the full original
\texttt{normUn\_gamma} statement, for the actual arithmetic group.

\begin{proof}
The exact \texttt{ModelBase.gam} condition says that $\gamma$
is unitary, preserves the actual $\mathcal O_E$-lattice, and
satisfies the half-level $l.M$ and whole-lattice $l.A$ congruences.
It therefore induces
$a=\phi^{-1}\gamma\phi:\Lambda\simeq_{\mathcal O_E}\Lambda$
with precisely the assumptions of NFL4.  The period formulas
give
\[
 P_{\gamma x}(av)=P_x(v).
 \tag{NFL.21}
\]
For early columns this is the unitary identity for
$\operatorname{herm}_G(\gamma\operatorname{col}_sX,\gamma v)$;
for late columns it is $(\gamma X)^{-1}\gamma v=X^{-1}v$.
Thus NFL7 applies with $H_0=I$.  The base points and eigenframes
are literally the same.  Normalization in (NFL.15) gives
$t_0^2=1$.  Take every eigenframe entry of $t$ equal to one and
its section entry equal to $t_0$.  This is the required
\texttt{FrBor}; its character is $t_0^n=1$ because $n$ is even.
The argument keeps the full original level conditions; it does
not replace $\Gamma$ by a convenient subgroup.
\end{proof}
% END NFL10

\subsection{Algebraic evaluations and holomorphic readings}

The holomorphy claims NFL12--NFL16 use only their supplied
$F,M,T$ and the hypotheses stated there, together with the two
retained section laws where indicated.  They do not require
\texttt{GoodEmb}, \texttt{ModelBase} or \texttt{NegXi}; those
conditions enter the construction and transformation laws.

% BEGIN NFL11
\paragraph{NFL11: actual evaluation is well defined and has the required weight law.}
\label{claim:NFL11}
At the retained \texttt{SecBaseChangeSt} premise, the original
\texttt{EvalWD} holds: every natural rule on \texttt{TRing}s
has the same value through any two test-ring presentations of
one framed point, over \emph{any} coefficient algebra $S$.
Moreover, for every reduced $c.F$-algebra $S$, every $a,n$,
every $f\in\operatorname{sec}(X_o,c,a,n)$ and every actual
$\operatorname{FrBor}(t,y,y')$ over $S$,
\[
 \operatorname{evR}(f,y')
   =t(\operatorname{none})^n
       \prod_{\rho,i}t(\operatorname{some}(\rho,i))^{a(\rho,i)}
          \operatorname{evR}(f,y).
 \tag{NFL.22}
\]
No reducedness is required in the first assertion, and no
weight-domination or parity condition is required in the second.

\begin{proof}
TDR5--TDR6 give, on each affine chart of the actual moduli scheme,
the universal frame ring
\[
 A=R[\{g_{\rho,ij}\},s,
           (s\prod_\rho\det g_\rho)^{-1}]
 \tag{NFL.23}
\]
and its actual universal framed point.  This is a reduced
finite-type algebra and hence a literal \texttt{TRing} by TDR12.
Its classifier is unique over every target ring, since it records
the base point, each eigenframe entry and the section generator.
Fix a maximal ideal $\mathfrak m\subset S$ and a framed point $w$.
Choose a chart through its underlying point over $S_\mathfrak m$;
all frames and the section generator give a unique map
$A\to S_\mathfrak m$ classifying $w_\mathfrak m$.

If $w$ was presented by a test point over $p.R$, the preimage
of this chart contains the point induced by $\mathfrak m$.
A principal neighborhood $p.R_g$ suffices.  After further such
localization, the finitely many entries of the eigenframes and
the section generator give its classifier $A\to p.R_g$.
This is exactly the finite localization capture of TDR8 applied
to the universal frame chart; $p.R_g$ remains reduced and of
finite type.  Naturality of the rule identifies its value from
$p$ with the image of its value at the universal point of $A$.
For a second test presentation the classifier into $S_\mathfrak m$
is the same by uniqueness.  Thus the two values agree in every
$S_\mathfrak m$ and consequently in $S$: a nonzero difference
has an annihilator contained in some maximal ideal and remains
nonzero in that localization.  This proves the full
\texttt{EvalWD}, even for nonreduced $S$.  Notice that the
argument used local universal charts, not a potentially
nonreduced common quotient as a new test ring.

Over (NFL.23) adjoin the entries of lower triangular change
matrices in the original eigenframe indices, their diagonal
inverses, and a section-scaling unit and its inverse.  The resulting
ring $B$ is again reduced and of finite type.  It carries two
actual framed points and their universal \texttt{FrBor} relation.
The defining weight axiom of $f$ gives (NFL.22) in $B$.
Any specified triangular change over $S_\mathfrak m$ specializes
this universal change, including its diagonal and section units.
The preceding evaluation argument identifies its two evaluations
with those of $y,y'$.  Hence (NFL.22) holds locally at every
maximal ideal and therefore in $S$, whenever both points have
test presentations.  TDR13 supplies those presentations for
every reduced $S$.  All pullbacks here are the literal
\texttt{FrMap} and \texttt{SecBC} of TDR3; their composition
and cancellation follow from the two fiber-product projections
and the chart section equalities.
\end{proof}
% END NFL11

% BEGIN NFL12
\paragraph{NFL12: local algebraic coordinates for the actual period family.}
\label{claim:NFL12}
At each $b_0$ of the open period-coordinate domain there are a
neighborhood $B$, a literal test ring $R$, an actual framed point
$y_0$ over $R$, and algebra maps $\psi_b:R\to\mathbb C$ for
$b\in B$ such that the base change $z_b$ of $y_0$ lies above
the actual point $F.u(P_b)$, and every $b\mapsto\psi_b(r)$ is
holomorphic.  The ring and frame are fixed on $B$ before $b$
and before any natural rule are chosen.

\begin{proof}
HUF21 says $b\mapsto F.u(P_b)$ is holomorphic into the actual
moduli scheme over the fixed coefficient embedding $\sigma$.
Choose one of the affine moduli charts with eigenbasis and
theta generator furnished by TDR5 through its value at $b_0$.
The preimage is open; shrink to a neighborhood $B$ lying in it.
The morphism into this affine chart gives the unique evaluation
homomorphism $\psi_b:R\to\mathbb C$.  Every $r\in R$ is a
regular function on the chart and hence its composition with
the holomorphic period family is holomorphic.
Base change of its fixed framed point is legitimate for every
$\psi_b$ by TDR3, using the retained section base-change law.
It gives $z_b$ with the original point, dual-number eigenframe,
and section-chart values, not just an analytic trivialization.
The finite-type reduced ring is a \texttt{TRing} by TDR12.
\end{proof}
% END NFL12

% BEGIN NFL13
\paragraph{NFL13: the change of eigenframes is holomorphic.}
\label{claim:NFL13}
More generally let $R$ be any commutative coefficient algebra,
$y_0$ a framed point, and $\psi_b:R\to\mathbb C$ such that all
its scalar evaluations are holomorphic on an open $B$.
Let $z_b$ be its actual base change above $F.u(P_b)$, and let
$y_b$ be a theta point whose eigenframe is the derivative of
$\exp_b$ at the standard basis.  The unique eigenframe change
matrices in $\operatorname{FrChg}(G_b,a_b,z_b,y_b)$ are
holomorphic functions of $b$.  No assumption on the section
factor $a_b$ is needed for this conclusion.

\begin{proof}
Only finitely many coefficients from $R$ are involved near a
fixed $b_0$: the algebraic group is of finite presentation over
$R$, and the chosen regular functions and differentials have
finite polynomial expressions on local charts.  Choose the
finitely many chart denominators whose images at $b_0$ are
nonzero and shrink $B$ until their images remain nonzero.
Evaluation of each resulting fraction is holomorphic because
its numerator and denominator are evaluations of elements of
$R$.  No complex analytic space structure on
$\operatorname{Spec}R$ is asserted or needed.
Near any chosen base point, smoothness of the actual abelian
scheme gives $g$ regular functions whose relative differentials
are a basis at its zero section.  One may choose them by
selecting a nonzero minor in a finite presentation of the
relative cotangent module.  Shrink the neighborhood so that
this minor stays invertible.
Compose these functions with the HUF holomorphic exponential
$(b,z)\mapsto\exp_bz$.  Their derivatives in $z$ at zero form
a holomorphic invertible Jacobian matrix.  Holomorphy of the
derivatives follows from the Cauchy integral formula on a
fixed small polydisc after shrinking $B$.
Evaluating the same differentials on the algebraic frame of
$y_0$ gives elements of $R$; their images under $\psi_b$ are
holomorphic by hypothesis.  Multiplication by the inverse
Jacobian therefore expresses $z_b$'s tangent vectors in
exponential coordinates by a holomorphic invertible matrix.
Its entries outside each CM eigenspace are zero, since both
sets of vectors have the prescribed $\mathcal O_E$ action.
Inverting each block gives $G_b$ (or its transpose, according
to whether the frame vectors are listed as rows or columns).
In the source convention, the rows of $G_b$ are precisely the
coefficients of $y_b.u_i$ in the frame $z_b.u_j$.
Every entry is holomorphic, since its determinant is nowhere
zero.  This also proves that these matrices are the literal
dual-number combination matrices; TDR1 identifies their
coefficients with the differentials just evaluated.
\end{proof}
% END NFL13

% BEGIN NFL14
\paragraph{NFL14: the actual theta section scaling factor is holomorphic.}
\label{claim:NFL14}
In NFL13 suppose $y_b$ is an actual theta point with scalar
$\kappa(b)$ holomorphic on $B$.  Then the section factor $a_b$
in $\operatorname{FrChg}(G_b,a_b,z_b,y_b)$ is holomorphic.
This holds for the same arbitrary $R$ and all holomorphic
evaluation homomorphisms, with the original section charts.

\begin{proof}
Fix $b_0$.  The pulled-back section $s_{0,b_0}$ generates the
nonzero one-dimensional complex section space.  Hence it does
not vanish everywhere on the complex fiber: this follows
either from its actual local regular functions on the reduced
complex variety or from the complex-fiber section theorem.
Choose a point where its value is nonzero and an original
line chart containing that point.  Surjectivity of the actual
exponential gives a $w_0\in\mathbb C^g$ representing it.
For $b$ near $b_0$, $\exp_bw_0$ remains in the same chart.
Its lift to the family pulled back from $y_0$ is holomorphic:
in finite local charts its coordinates are the evaluations
$\psi_b$ and the HUF holomorphic exponential coordinates,
and their defining base equations agree.  More explicitly,
the fixed algebraic section on this chart is a finite
polynomial or rational expression in those coordinates and
finitely many coefficients from $R$.  Localize at its finitely
many denominators and shrink $B$ so that their evaluations
stay nonzero, as in NFL13.  This proves holomorphy directly
without analytifying the arbitrary ring $R$.
Thus the value $v(b)$ of the fixed algebraic section
$s_0$ at this point is holomorphic, and $v(b_0)\ne0$.
Shrink the neighborhood until $v$ is nowhere zero.

The theta section at that point has the holomorphic value
\[
 \kappa(b)\theta[\alpha,\beta](\tau_b,h_bw_0)fr_j(b,w_0).
\]
Normal convergence from NFR5 gives holomorphy of the theta
factor; the other factors are the actual holomorphic data.
The chart equality in \texttt{SecBC} says this value is
$a_bv(b)$.  Dividing by $v(b)$ proves local holomorphy of $a_b$.
Since $b_0$ was arbitrary, it proves the assertion on $B$.
The chosen evaluation point is allowed to depend on $b_0$;
there is no assumption that theta is nonzero at the zero section.
\end{proof}
% END NFL14

% BEGIN NFL15
\paragraph{NFL15: natural rules read holomorphically on algebraic frame charts.}
\label{claim:NFL15}
More generally, let $V$ be any complex normed vector space,
$B\subset V$ any set, $R$ any literal test ring, $y_0$ an
actual framed point over $R$, and $\psi_b:R\to\mathbb C$
evaluation homomorphisms with actual base changes $z_b$.
Suppose every scalar evaluation is differentiable on $B$ in
the source's within-set sense, and suppose an actual
framed point $y_b$ is related to $z_b$ by
$\operatorname{FrChg}(G_b,a_b,z_b,y_b)$, with every entry of
$G_b$ and $a_b$ differentiable on $B$.  Then for every natural
rule $f$ on the original test rings,
$b\mapsto\operatorname{evR}(f,y_b)$ is differentiable on $B$.
In the open finite-dimensional setting of NFL12 this is
holomorphy.  No weight condition is necessary here.

\begin{proof}
Use the universal full frame ring (NFL.23) over this $R$.
Its base-changed frame with matrices $G_b$ and section factor
$a_b$ is classified by one map $\eta_b:A\to\mathbb C$.
The entries of $G_b$ and $a_b$ are differentiable on $B$ by hypothesis;
their determinants and $a_b$ are nonzero because both sides
are actual bases and section generators.  Their inverses
are therefore differentiable on $B$.  Every element of $A$ is a
polynomial in the entries, the base elements of $R$, and these
inverses.  Consequently every $b\mapsto\eta_b(v)$ is differentiable on $B$.
Let $y_A$ be its universal framed point.  NFL11 and naturality
give the literal evaluation identity
\[
 \operatorname{evR}(f,y_b)=\eta_b\bigl(f_A(y_A)\bigr).
\]
The element on the right inside parentheses is fixed before
$b$ varies.  This proves the assertion.  It is the universal
frame representing property, including the section factor,
that makes the result hold for every rule simultaneously.
\end{proof}
% END NFL15

% BEGIN NFL16
\paragraph{NFL16: all readings of the normalized points are holomorphic.}
\label{claim:NFL16}
For every $a,n$ with $\operatorname{Dom}(a,n)$, every
$f\in\operatorname{sec}(X_o,c,a,n)$ and every component $k$,
the original reading $\operatorname{rdU}(f,k)$ is holomorphic
on $(M.X_k).\mathrm{dom}$.  This is the complete original
\texttt{differentiableOn\_rdU} conclusion for any family satisfying
\texttt{IsNormUn}, with its unrestricted pointwise square roots.

\begin{proof}
At each period coordinate $b$, take the actual theta point
$y_b^1$ with scalar $1$, whose existence is NFR6.
On a neighborhood from NFL12 compare it with the base change
$z_b$ of the fixed algebraic framed point.
The eigenframe change is holomorphic by NFL13 and the section
factor by NFL14, since the scalar $1$ is holomorphic.
NFL15 therefore proves that
$b\mapsto\operatorname{evR}(f,y_b^1)$ is holomorphic.

The normalized point $u_k(T.\mathrm{co}(b))$ has the same
base point and eigenframes as $y_b^1$; its actual section is
$\kappa_b$ times the section of $y_b^1$.  This is the
\texttt{FrBor} with eigenframe diagonal entries $1$ and section
entry $\kappa_b$.  Apply NFL11 over $\mathbb C$ and write $n=2m$:
\[
 \begin{split}
 \operatorname{evR}(f,u_k(T.\mathrm{co}(b)))
   &=\kappa_b^n\operatorname{evR}(f,y_b^1)\\
   &=(C\det h_b)^m\operatorname{evR}(f,y_b^1).
 \end{split}
 \tag{NFL.24}
\]
The last expression is holomorphic and is independent of the
chosen square root at every point.  The coordinate map
$T.\mathrm{co}$ and its inverse are complex linear; transferring
back to the open frame domain proves the assertion.
The source extends $\operatorname{rdU}$ by zero outside the
domain; this has no bearing on holomorphy within the stated
domain and no boundary regularity is claimed.
\end{proof}
% END NFL16

% BEGIN NFL17
\paragraph{NFL17: one family satisfies every field of the original FramesOK.}
\label{claim:NFL17}
For the original supplied $F,M$ and all their original data,
there is one family
\[
 u:\prod_k(M.X_k).\mathrm{dom}\longrightarrow
       \operatorname{FrPt}(X_o,c,\mathbb C_\sigma)
\]
satisfying the complete original \texttt{FramesOK} record.
Every point is also \texttt{FromTest}.
The family is chosen before the components, frames, arithmetic
elements, triangular moves, weights and natural rules that occur
in its universal assertions.

\begin{proof}
Choose exactly the family of NFR13, with one theta datum per
component and the pointwise square roots of NFR6.
NFR12 supplies its first field, the full \texttt{FrameOK} at
every point.  NFL9 supplies the second field, with its single
triangular change witness valid for all allowed weights.
NFL10 supplies its third field, the arithmetic invariance
witness valid for all allowed weights.  NFL16 supplies its
fourth field, holomorphic readings for every allowed section.
These are the four fields in the literal definition, with no
replacement of the chosen family at any join.
TDR13 gives \texttt{FromTest} at every point because
$\mathbb C_\sigma$ is reduced.  The only additional law premises
are precisely the section base-change and twist laws stated
at the beginning of this chapter.
\end{proof}
% END NFL17

\subsection{Recovering frames and identifying their fibers}

% BEGIN NFL18
\paragraph{NFL18: every original period map is a literal frame period map.}
\label{claim:NFL18}
For every $P:\Lambda\to\mathbb C^g$ satisfying the original
\texttt{IsPeriod} condition there is $x\in\mathcal D$ whose
literal NFR1 period map is $P$.  Together with NFR2 this gives
the entire original \texttt{exists\_framePeriods}, including
its previously unused surjectivity clause.

\begin{proof}
The lattice and its prescribed model identify
$\Lambda\otimes\mathbb R$ with
$\bigoplus_{\tau\in\Psi}\mathbb C^r$.
The full lattice condition in \texttt{IsPeriod} makes the real
extension of $P$ an isomorphism.  The $\mathcal O_E$ eigenvalue
condition extends first to $E$ and then to $E\otimes\mathbb R$.
In the unreversed column dictionary, set
\[
 u_{\tau,s}=\begin{cases}
   \overline{P_{\tau,s}},&s<q,\\ P_{\tau,s},&s\ge q.
   \end{cases}
\]
Each block $u_\tau$ is $\tau$-linear in the model coordinate
$v_\tau$, so $u_\tau=M_\tau v_\tau$ for an invertible complex
matrix $M_\tau$.  Put $Y_\tau=M_\tau^{-1}$ and
$S_\tau=Y_\tau^tG_\tau\overline{Y_\tau}$.

Let $K$ be the positive definite Hermitian form provided by
\texttt{IsPeriod}, so $\operatorname{Im}K(Pv,Pw)=E_\Lambda(v,w)$.
The identity $E_\Lambda(av,w)=E_\Lambda(v,\bar a w)$, first
on the lattice and then over $\mathbb R$, implies
$K(az,w)=K(z,\bar a w)$: their imaginary parts agree, and
replacing $z$ by $iz$ identifies their real parts.
Thus distinct eigenvalue types are orthogonal for $K$.
In particular its early and late blocks at each $\tau$ are
separate, since their types are $\bar\tau$ and $\tau$.
Write these positive definite blocks as $K_+,K_-$, and put
$a_\tau=\operatorname{Im}\tau(-\xi_0)>0$.
The trace identity (NFL.1), on vectors in this real factor, is
$E_\Lambda=-2a_\tau\operatorname{Im}\operatorname{herm}_{G_\tau}$.
In the $u$ coordinates, comparison of imaginary parts gives
\[
 S_\tau=
 \begin{pmatrix}
  \overline{K_+}/(2a_\tau)&0\\
  0&-K_-/(2a_\tau)
 \end{pmatrix}.
 \tag{NFL.25}
\]
The imaginary part determines a first-linear Hermitian form,
since replacing one input by $i$ recovers its real part; hence
this is an equality of matrices, not just of their imaginary
parts on the diagonal.

Let $S_+$ be its early block and define
\[
 B_\tau=\operatorname{diag}((S_+^{-1})^t,I_{r-q}),
 \qquad X_\tau=Y_\tau B_\tau.
 \tag{NFL.26}
\]
For $s<q$, the $s$th row of $B_\tau^tS_\tau$ is the $s$th
standard row, so
\[
 \operatorname{herm}_{G_\tau}(\operatorname{col}_sX_\tau,v_\tau)
      =\overline{u_{\tau,s}}=P_{\tau,s}.
\]
For $s\ge q$, $(X_\tau^{-1}v_\tau)_s=u_{\tau,s}=P_{\tau,s}$.
These are exactly the two original period formulas, including
the reversal of the early indices.  The Gram matrix of $X_\tau$
is $B_\tau^tS_\tau\overline{B_\tau}$, with positive definite
early block and negative definite late block.  Every leading
minor consequently has the prescribed sign: before $q$ it
is positive, and after $q$ it is $(-1)^{s-q}$ times a positive
number.  Hence $X$ lies in the actual domain $\mathcal D$.
Empty blocks use the empty identity matrix and determinant one,
so the argument includes $q=0$ and $q=r$.
\end{proof}
% END NFL18

% BEGIN NFL19
\paragraph{NFL19: the full original triangular covering theorem.}
\label{claim:NFL19}
For every actual complex framed point $y$, there are $k$,
$x\in(M.X_k).\mathrm{dom}$ and the literal period map $P_x$
such that $F.u_k(P_x)=y.x$ and the standard exponential
derivatives equal every prescribed vector $y.u_{\rho,i}$.
Consequently the original \texttt{CoverTri} holds, with unit
diagonal coefficients and zero off-diagonal coefficients.
This statement concerns every derivative realization in the
source's quantifier order, not just a chosen realization.

\begin{proof}
The onto field of HUF21 supplies a component $k$ and an admissible
period map $P$ above $y.x$.  Its exponential derivative is a
complex-linear isomorphism onto the actual Lie algebra.
Pull back each $y.u_j$ to a vector $w_j\in\mathbb C^g$.
The actual eigenbasis property of $y$ implies that these vectors
form a basis and that their columns have support only in the
corresponding CM eigenspace.  Thus the matrix $W$ with columns
$w_j$ is invertible and block diagonal by original embedding.
Put $V=W^{-1}$ and $P'=VP$.  The image is again a full lattice,
the types are unchanged, and a witnessing positive Hermitian
form is $K'(z,w)=K(V^{-1}z,V^{-1}w)$.  Hence $P'$ satisfies
the exact \texttt{IsPeriod} condition.

Use the identity lattice isometry and the linear change $V$ in
the HUF family identification.  Its level conditions are
trivial, and it gives
$F.u(P')=F.u(P)=y.x$ and $\exp_{P'}(Vz)=\exp_P(z)$.
Since $Vw_j=e_j$, the standard derivative for $P'$ is exactly
$y.u_j$.  NFL18 now supplies a frame $x$ with $P_x=P'$.
For any other vectors $v_{\rho,i}$ satisfying the source's
specified standard derivative relation, uniqueness of the
dual-number derivative gives $v_{\rho,i}=y.u_{\rho,i}$.
Therefore the coefficient vector $cf_j=\delta_{ij}$ witnesses
every required \texttt{IsCombAt}; its diagonal is $1$ and
the higher coefficients vanish.  This proves the full
\texttt{CoverTri} quantifiers without using the section of $y$.
\end{proof}
% END NFL19

% BEGIN NFL20
\paragraph{NFL20: the full original triangular fiber theorem.}
\label{claim:NFL20}
Take arbitrary components $k,k'$, frames $x,x'$ in their actual
domains, their literal period maps $P,P'$, and any tangent
families $v,v'$ realizing the standard exponential derivatives.
Suppose the base points are equal and each $v'_{\rho,i}$ is a
linear combination of the $v_{\rho,j}$ with $j\le i$, with
unit coefficient at $i$.  Then there exist
$\gamma\in(M.X_k).\Gamma$ and $b\in\operatorname{Bor}(\Psi,r)$
such that
\[
 (k',x')=(k,\gamma x b)
 \tag{NFL.27}
\]
as an equality in the original dependent union of component
frame spaces.  This is the full original \texttt{FibTri}.

\begin{proof}
NFR2 makes $P,P'$ admissible.  The disjointness field of HUF21
first gives $k'=k$; its fiber field then gives an actual
$\mathcal O_E$-lattice isometry $a$ preserving $E_\Lambda$,
the two original level conditions, and a complex-linear
isomorphism $H_0$ with $P'a=H_0P$.
Extend $a$ to the model $E$-space via $\phi$.  This gives
$\gamma\in\operatorname{GL}_r(E)$ on the same full lattice.
As in NFL4, trace nondegeneracy proves that $\gamma$ is unitary.
The two level conditions transfer through $\phi$ unchanged,
so the exact \texttt{ModelBase.gam} equivalence gives
$\gamma\in(M.X_k).\Gamma$.  In particular this is a witness
in the original group, not just in its rational unitary group.

By NFL10's period computation, $P_{\gamma x}a=P_x$.
Since $a$ is surjective, $P_{x'}=H_0P_{\gamma x}$.
The family identification gives
$\exp_{x'}(H_0z)=\exp_xz$ on the common fiber.
Let $C$ have as its $i$th column the coefficients expressing
$v'_i$ in the vectors $v_j$; its blocks are upper triangular
in the original indices.  Differentiating the last exponential
identity and using injectivity of the derivative gives
$H_0C=I$.  Hence $H_0=C^{-1}$ has the same block triangular
support, including the separation by original embedding.

Put $b_\tau=(\gamma_\tau X_\tau)^{-1}X'_\tau$.
We explain why it is upper triangular, rather than merely
invertible.  Equality of period maps on the full lattice
extends to every real model vector.  Comparing the early
period coordinates expresses each early column of $X'$ as a
linear combination of the early columns of $\gamma X$:
nondegeneracy of $G_\tau$ permits cancellation against all
vectors.  Thus the lower-left block of $b_\tau$ is zero.
The early-to-early formula is $H_{0,++}=b_{++}^t$ in unreversed
column coordinates.  Because those indices are reversed in
the original eigenframe order, the triangular support of
$H_0$ says $b_{++}$ is upper triangular.  The late formula is
$H_{0,--}=b_{--}^{-1}$; its indices retain their order, so
$b_{--}^{-1}$ and therefore $b_{--}$ are upper triangular.
Both diagonal blocks are invertible.  Together with the
vanishing lower-left block this makes the whole $b_\tau$
upper triangular and invertible.  Its possible upper-right
block is allowed.  Thus $b\in\operatorname{Bor}$ and (NFL.27)
holds.  All steps used only the supplied derivative relations,
so the conclusion is valid for the arbitrary $v,v'$ of the
original statement.
\end{proof}
% END NFL20

\subsection{The complete uniformization in the retained original scope}

% BEGIN NFL21
\paragraph{NFL21: the actual family is a full complex uniformization.}
\label{claim:NFL21}
Let $U=\operatorname{unfOf}(e,F,M,u)$ for any $u$ satisfying
the complete original \texttt{FramesOK}.  At the original
\texttt{GoodEmb}, \texttt{ModelBase}, \texttt{NegXi} and
\texttt{SecBaseChangeSt} premises, $U$ satisfies the full
\texttt{IsUnf}.  In particular this applies to the single
family of NFL17.  The additional twist law is needed for
that construction, not as an additional premise to this join.

\begin{proof}
We verify the clauses, including the data that accompanies the
orbit assertions.  The rank $r=D_w.p+D_w.q$, the value
$(U.X_k).p_0=D_w.q$, and the exact host CM type are the rank
and type clauses of \texttt{ModelBase}.
For an arbitrary framed point $y$, NFL19 supplies $k,x$ whose
base point and standard derivatives are $y$'s.  The
\texttt{FrameOK} field for $u_k(x)$ has the same literal
period map and derivatives.  Its section and $y.s$ are both
generators of the same line, so their ratio is a unit.
Together with the identity eigenframe coefficients this gives
$\operatorname{FrBor}(t,u_k(x),y)$.  This is the covering clause.

If $\operatorname{FrBor}(t,u_k(x),u_{k'}(x'))$, its base equality
and triangular derivative clauses meet NFL20 exactly, giving
$(k',x')=(k,\gamma xb)$ for original $\gamma,b$.
Conversely, for such $\gamma,b$, the arithmetic and Borel
fields of \texttt{FramesOK} give successive actual
\texttt{FrBor} relations.  Multiplying their triangular
matrices and their section scalars gives their composite
\texttt{FrBor}.  This proves both directions of the exact
dependent-component fiber equivalence.

For $f\in\operatorname{sec}(a,n)$ with $\operatorname{Dom}(a,n)$,
its reading is holomorphic by the fourth \texttt{FramesOK}
field.  The domain is stable under both actions: unitary
left multiplication preserves its Gram matrices, and upper
triangular right multiplication multiplies each signed leading
minor by the positive number $|\det b_{[0,s)}|^2$.
Apply NFL11 to the two \texttt{FrBor} relations of the second
and third fields.  Their character equalities give
\[
 \operatorname{rdU}(f,k)(xb)=bchar(\operatorname{laU}(a,n,k))(b)
                              \operatorname{rdU}(f,k)(x),
 \qquad
 \operatorname{rdU}(f,k)(\gamma x)=\operatorname{rdU}(f,k)(x).
 \tag{NFL.28}
\]
The latter is precisely the required arithmetic character law,
because \texttt{ModelBase.chi} makes that character identically
one for every weight.  These are all clauses of the original
$\operatorname{IsSec}$ statement.

Finally take the witness $g,e,F$, component map $k\mapsto k$,
and the very same maps $M.\phi_k$ in the last existential field
of \texttt{IsUnf}.  Its trace identity is (NFL.1).  Each
\texttt{FrameOK} supplies the literal period formula,
\texttt{IsPeriod}, equal base point and all standard derivative
identities.  At every frame with unit absolute diagonal minors,
it also supplies the positive Hermitian matrix, actual
holomorphic nowhere-zero local frames and their cocycle and
rigidification, the unitary lattice multiplier, the actual
section-value identity, the additive lattice fundamental
domain and the Gaussian mean exactly one.  These are exactly
the remaining conjuncts of \texttt{IsUnf}, already proved in
NFR12, with the same section and same period lattice.  Thus no
existential witness has been replaced at the last join.
\end{proof}
% END NFL21

% BEGIN NFL22
\paragraph{NFL22: principal level and every cusp for the prescribed Hermitian space.}
\label{claim:NFL22}
Fix arbitrary original $E,H,D_w,l,X_o,c,\sigma,r,\xi_0,G$ and
$\phi_0:D_w.\mathrm{La}\to_{\mathcal O_E}E^r$ with
$1\le D_w.p+D_w.q$, \texttt{GoodEmb}, $\bar\xi_0=-\xi_0$,
$\operatorname{Im}\tau(\xi_0)<0$ on the host type,
$\phi_0$ injective with full $E$-span, $G$ Hermitian and invertible,
and
\[
 D_w.E_F(v,w)=\operatorname{Tr}_{E/\mathbb Q}
       (-\xi_0\operatorname{herm}_G(\phi_0v,\phi_0w)).
\]
At the supplied \texttt{Moduli}, the exact retained
\texttt{SecBaseChangeSt} and \texttt{SecTwistSt}, and the
inherited HUF/PUN classical-foundation scope, there exists one
$U:\operatorname{Unf}(X_o,c,\sigma,r)$ with all the following:
\[
 \begin{gathered}
 \operatorname{IsUnf}(U),\qquad
 \operatorname{HasPrinc}(U,l.M'),\qquad U.\mathrm{Comp}\ne\varnothing,\\
 \forall k,\lambda,\gamma,\quad (U.X_k).\chi(\lambda,\gamma)=1,\qquad
 \forall k,\quad (U.X_k).G_m=G,\ (U.X_k).\xi_0=\xi_0,\\
 \forall k,x,\quad\operatorname{FromTest}(U.u_k(x)),\\
 \forall\rho,\quad\operatorname{FJ.LatZ}(U.X,\rho),\qquad
 l.A=1\Longrightarrow\operatorname{FJ.ExpoMatch}(U.X,1,\rho,l.M).
 \end{gathered}
 \tag{NFL.29}
\]
Here $l.M'=\operatorname{lcm}(l.M,l.A)$ is the literal original
modulus.  Every cusp $\rho$ and every frame point uses this
same $U$, chosen beforehand.

\begin{proof}
HUF21 supplies the full actual family $F$ for this moduli
problem; its component count is positive by the stipulated
positive rank.  PUN11 supplies $M$ with the exact
\texttt{ModelBase} and \texttt{NegXi}, with $G$ and $\xi_0$
equal to the prescribed values on every component.  NFL17
supplies one full \texttt{FramesOK} family and NFL21 gives
$\operatorname{IsUnf}$ for its $U$.

To verify the literal \texttt{HasPrinc}, take a unitary lattice
stabilizer $\gamma$ congruent to $I$ on the whole lattice
modulo $N=l.M'$.  Since $l.M\mid N$, its difference from
the identity lies in $l.M\Lambda$ for every lattice vector,
in particular on the host half-ideal sublattice.  Since
$l.A\mid N$, the second required congruence also holds.
These implications mean actual vectors: if
$(\gamma-I)v=Nw$ and $N=md$, use the vector $dw$ as the
$m$-divisibility witness.  The exact \texttt{ModelBase.gam}
equivalence therefore puts $\gamma$ in $\Gamma$.
This proves \texttt{HasPrinc} without imposing any extra
congruence on $U$.  PUN6 in fact also identifies the original
group with the principal subgroup at this modulus.

Nonempty components and trivial characters are the already
specified fields of $F$ and $M$.  The equality of Hermitian
data is the prescribed-space conclusion of PUN11.  TDR13
gives the literal \texttt{FromTest} witnesses for all its
complex framed points.  PUN12--PUN14 give the full center
lattice for every cusp of these actual arithmetic groups;
PUN15--PUN19 identify the original integral Fourier exponent
lattice when $l.A=1$, producing the asserted
\texttt{ExpoMatch} at level $l.M$.  Their statements are
uniform in the cusp and use these same lattices and groups,
so they may be joined after $U$ is chosen.

This reconstructs the mathematical return of the source's
\texttt{exists\_unf\_hasPrinc} at the two explicitly retained
section laws.  The source obtains those laws from its global
law packages; their global suppliers have not been proved in
this packet.  Accordingly (NFL.29) is a complete theorem
with these stated premises, not an assertion that the
unconditional source wrapper or a selected whole obligation
has been independently reconstructed.
\end{proof}
% END NFL22

\subsection{Source and inherited-proof map}

\begingroup\raggedright
The metadata records every complete pinned file, exact range,
SHA256, and the proof-block hashes of all claims and inherited
inputs.  The following identifies the mathematical source of
each group.  Paths in this paragraph are relative to
\texttt{src/\\allowbreak{}Development/\\allowbreak{}Abelian\allowbreak{}Variety\allowbreak{}Theory/\\allowbreak{}Proofs/\\allowbreak{}}.
NFL1--NFL2 reconstruct the finite Gauss sum and Poisson argument
in \texttt{HomogeneousIdeals/}\\
\texttt{From\allowbreak{}Blocks\allowbreak{}Map\allowbreak{}Int\allowbreak{}Cast\allowbreak{}Mem\allowbreak{}Symplectic\allowbreak{}Group.lean}, together
with its finite-form diagonalization supplier.  NFL3 reconstructs
\texttt{Unitary\allowbreak{}Existence/\\allowbreak{}Two\allowbreak{}Dvd\allowbreak{}Multiplier\allowbreak{}Of\allowbreak{}Theta\allowbreak{}ZAct\allowbreak{}Eq\allowbreak{}Unit\allowbreak{}Mul.lean}
and \texttt{RatioSqEqDetJAut.lean}.
NFL4--NFL7 correspond to \texttt{Cm\allowbreak{}Unit\allowbreak{}Pow\allowbreak{}Four\allowbreak{}Idempotent\allowbreak{}Mod\allowbreak{}Eight.lean},
\texttt{ChangeOfNormalisation.lean},
\texttt{ExistsSecGeneratorFibre.lean},
\texttt{Theta\allowbreak{}Frame\allowbreak{}Det\allowbreak{}Mul\allowbreak{}Sq\allowbreak{}Invariant.lean} and
\texttt{Theta\allowbreak{}Points\allowbreak{}Change\allowbreak{}Level\allowbreak{}Isometry\allowbreak{}Linear\allowbreak{}Map.lean} in that
same unitary-existence directory.

NFL8--NFL10 reconstruct \texttt{Triangular\allowbreak{}Frame\allowbreak{}Moves\allowbreak{}Period\allowbreak{}Block\allowbreak{}Triangular.lean},
\texttt{Normalised\allowbreak{}Framed\allowbreak{}Point\allowbreak{}Triangular\allowbreak{}Frame\allowbreak{}Weight\allowbreak{}Rule.lean},
\texttt{PerFrLact.lean} and
\texttt{Normalised\allowbreak{}Framed\allowbreak{}Point\allowbreak{}Arithmetic\allowbreak{}Group\allowbreak{}Invariant\allowbreak{}Sig.lean}.
NFL11 is the evaluation and weight portion of
\texttt{Framed\allowbreak{}Moduli\allowbreak{}Algebraic\allowbreak{}Laws\allowbreak{}Chart\allowbreak{}Weight\allowbreak{}Integrality.lean},
using the actual universal charts already proved in TDR;
it makes no claim here to that file's separate integral-lattice
or unit-independence conjuncts.
NFL12--NFL16 reconstruct, respectively,
\texttt{Exists\allowbreak{}Chart\allowbreak{}Holomorphic\allowbreak{}Coordinates\allowbreak{}Near\allowbreak{}Frame.lean},
\texttt{Differentiable\allowbreak{}On\allowbreak{}Exponential\allowbreak{}Frame\allowbreak{}Matrix.lean},
\texttt{Differentiable\allowbreak{}On\allowbreak{}Theta\allowbreak{}Algebraic\allowbreak{}Generator\allowbreak{}Factor.lean},
\texttt{Differentiable\allowbreak{}On\allowbreak{}Natural\allowbreak{}Rule\allowbreak{}Holomorphic\allowbreak{}Frames.lean} and
\texttt{Differentiable\allowbreak{}On\allowbreak{}Readings\allowbreak{}Normalised\allowbreak{}Theta\allowbreak{}Points.lean}.

NFL18--NFL20 reconstruct the onto clause of
\texttt{ExistsFramePeriods.lean} through
\texttt{ExistsFrameOfIsPeriod.lean}, then
\texttt{ExistsInverseBlock.lean} and
\texttt{BlockTriangularOfMulEqOne.lean}, with their exact period
and derivative identifications.
NFL17 and NFL21--NFL22 perform the original joins of
\texttt{Exists\allowbreak{}Framed\allowbreak{}Points\allowbreak{}All\allowbreak{}Normalisations\allowbreak{}Family\allowbreak{}Unif.lean},
\texttt{Family\allowbreak{}Uniformisation\allowbreak{}Is\allowbreak{}Complex\allowbreak{}Uniformisation.lean} and
\texttt{CongrOfDvdModulus.lean}.  The exact source
\texttt{NormConst}, \texttt{UnivFrame}, \texttt{IsCombAt},
\texttt{AlgebraMapCAt}, \texttt{UnFree} and \texttt{BaseNorm}
definitions are included in the manifest, so the source names
in these endpoints refer to their original records and
quantifier order.

The inherited human proofs are NFR1--NFR13, HUF1 and HUF21,
TDR1, TDR3, TDR5--TDR6, TDR8, TDR12--TDR13, PUN6,
PUN11--PUN19.  Each is bound to its exact earlier proof block
and keeps its earlier premises.  The present chapter has no
new Fourier-measure normalization, replacement section, hidden
holomorphic square-root choice, or altered coefficient field.
\par\endgroup
